
\lesson{27/10}
Since the off-diagonal elements represent transition apmlitudes, 
at equilibrium the matrix should be symmetric, i.e.
$$\rho_{\alpha\beta}=\rho_{\beta\alpha}.$$

The deistic matrix $\rho$ has the same role as $\rho(p,q)$ in the classical
pictures and obeys the equation
$$i\hbar\pdv[]{\rho}{t} = \left[H,\rho\right]\to \rho(t)=e^{-\frac{i}{\hbar}Ht}
\rho(0)e^{\frac{i}{\hbar}Ht}$$
Again, at equilibrium
$$\pdv[]{\rho}{t}=o\iff\left[H,\rho\right]=0$$
which means that $\rho$ is a function only of operators that commute with $H$:
the assumption taken is 
$$\rho = \rho(H)$$
and at most of the number operator $\hat N$ if the system is open.

At equilibrium a convenient base is that of the stationary states of $H$
$$[\rho,H]=0\to \rho = \sum_{n}^{}\rho_n\ket{n}\bra{n}$$

Recall that the Von Neumann Entropy is defined as
$$S=-\Tr(\rho\ln(\rho)) = -\sum_{i}^{}\lambda_i\ln\lambda_i$$
which measures the information stored in a state: $S=0$ for pure states,
$S>0$ for mixed states.
\section{Gibbs Statistical Ensembles}
Recalling that in classical SM
$$O_\text{obs}=\sum_{\nu}^{}P(\nu)O(\nu)=\langle O\rangle$$
where $\nu$ is the set of microstates of the system compatible with the
chosen macrostate. The goal of this part is to
\begin{itemize}
    \item Determine the general form of $P(\nu)$ for isolated, closed... systems.
    \item Express TD quantities in terms of $P(\nu)$.
\end{itemize}
\subsection{Microcanonical Ensemble}
It describes \textbf{isolated systems at fixed} $\mathbf{V,N,E}$: note that
there can be uncertainty on the value of $E$, but under the condition
$$\frac{\Delta E}{E}<<1.$$
At equilibrium, $P(\nu) = f(E)$ is constant, because $E(\nu)$ is, so recalling
$$\sum_{\nu=1}^{\Omega} P(\nu) = P\sum_{\nu=1}^{\Omega}=1 \iff
P\equiv\frac{1}{\Omega}$$
where $\Omega$ is the total number of microstates defined by E,V,N: this can be
interpreted as the fact that the system explored all microstate in the nonzero time
of the macroscopic measurement; this is called \textit{ergodicity}.

Although it is realized for almost all macroscopic systems, it is problematic
to establish rigorously [qualcosa su KAM theorem] and is surprisingly difficult to
achieve. $P(\nu)=\frac{1}{\Omega}$ is admissible if and only if the return time
of the system (i.e. the time for the system to return arbitrarily close to its
initial position) is comparable to the measurement time: the problem is that
the return time is typically very large. [qualcosa] the Kac Lemma
$$\tau_\text{ret} \propto\tau e^{\alpha N}$$
if $P(\nu)\propto e^{-\alpha N}$. [small things]\\

To find the TD quantities' expressions, $S(E,V,N)$ is needed:
consider a system subdivided into 2 isolated parts $1,2$; then the subdivision
becomes so that they can exchange only heat: the total energy is the same, but the
equilibrium variables vary. The number of microstates are
$$\Omega_1(E_1,V_1,N_1),\; \Omega_2(E_2,V_2,N_2)\to\Omega(E)=\Omega_1\Omega_2$$
Before (?) the second equilibrium, the situation is
$$\Omega(E) = \int \dd E_1 \Omega_1(E_1)\Omega_2(E-E_1)$$
then it reaches a situation in which the number of microstates is maximum,
with energies $E_1=\bar E_1,\;E_2=\bar E_2$.

$\Omega_1$ at equilibrium must be sharply peaked at the value $\bar E_1$, which
maximises the integrand:
$$0=\pdv[]{}{E_1} \left(\Omega_1(E_1)\Omega_2(E-E_1)\right) =
\pdv[]{\Omega_1}{E_1}(E_1)\Omega_2(E-E_1) - \Omega_1(E_1)\pdv[]{\Omega_2}{E_2}$$
$$\frac{1}{\Omega_1(E_1)}\pdv[]{\Omega_1}{E_1} = \frac{1}{\Omega_2(E_2)}
\pdv[]{\Omega_2}{E_2}$$
$$\pdv[]{\ln\Omega_1}{E_1} = \pdv[]{\ln\Omega_2}{E_2}$$
since the TD equilibrium condition is $T_1 = T_2$, which means
$$\frac{1}{T_1} = \frac{1}{T_2} \iff \eval{\pdv[]{S}{E_1}}_{V_1,N_1} =
\eval{\pdv[]{S}{E_2}}_{V_2,N_2}$$ 
comparing the two expressions yields
$$S_i\propto\ln\Omega_i \to S = k\ln\Omega$$
where $k$ can be experimentally proven to be a universal constant, the Boltzmann
constant
$$k = \frac{R}{N_A} = 1.38\cdot10^{-23}\text{ JK}^{-1}.$$
Since this arose from a confront of derivatives, a constant offset could have
been added, but $\Omega=1\to S=0$ is convenient, since Nernst's principle
prescribes that at $T=0$ the entropy must go to 0.

\textbf{Note} This form of $S$ preserves the fact that it is an extensive
quantity.

The general expression for entropy will be
$$S = -k\sum_{\nu}^{}P(\nu)\ln P(\nu)$$
\subsection{Example: Classical ideal gas}
The Hamiltonian is the free $N$ particle Hamiltonian
$$H = \frac{1}{2m}\sum_{i=1}^{N}p_i^2$$
and the microcanonical ensemble is a region of the phase space such as
$$E\leq H\leq E+\Delta E,\;\frac{\Delta E}{E}<<1$$
the possible values of the momenta live in a shell of width $\Delta E$ around
the hypersurface
$$\sum_{i=1}^{N}p^2_i = 2mE:=R^2(E)$$
which is a $(3N-1)$-sphere, so the volume of this region is, at first order, the
area of this sphere multiplied by its width:
$$V\approx \Sigma_{3N+1}(R(E))\cdot(R(E+\Delta E) - R(E))$$
\textbf{Note} The surface of a unit $D$-sphere.

Consider the integral
$$I_D = \int \dd^Dx e^{-(\mathbf{x})^2}=\left(\int \dd x e^{-x^2}\right)^D=
(\pi)^{D/2}$$
this can also be computed in (hyper-)spherical coordinates
$$I_D = \int \dd\Omega\int \dd r r^{D-1}e^{-r^2} $$
the integral over the hyper-solid angle $\Omega$ is the area of the
$(D-1)$-sphere! So
$$I_D = \Sigma_{D-1}\int_0^\infty \dd r r^{D-1}e^{-r^2}$$
letting $t=r^2\to \dd t = 2r\dd r$ yields
$$I_D = \Sigma_{D-1}\frac{1}{2}\int \dd t t^{\flatfrac{D}{2}-1}e^{-t} =
\Sigma_{D-1}\frac{\Gamma(\flatfrac{D}{2})}{2}$$
so
$$\Sigma_{D-1} = \frac{2\pi^{\flatfrac{D}{2}}}{\Gamma(\flatfrac{D}{2})}$$

